\section{Methodology}
\label{sec:method}

We investigate the detectability and latent separation of text origin modes through a two-stage experimental design.
First, we construct a controlled composite dataset.
Second, we evaluate zero-shot classification performance and embedding cluster separation.

\subsection{Dataset Preparation}
\label{subsec:data}

To isolate the 'mode' variable while controlling for content, we generated a balanced composite dataset of 150 samples (50 per mode). We utilized the following sources:
\begin{itemize}[leftmargin=*]
    \item \aigen: Samples were drawn from the \textsc{jjz5463} LLM-detection dataset \cite{jjz2024llm}, representing clean, model-generated text.
    \item \dictated: We simulated dictation artifacts by injecting phonetic homophone substitutions (e.g., "their" for "there") into source human texts, drawing from the \textsc{youngwoo3283} ASR dataset patterns \cite{youngwoo2024asr}.
    \item \typed: We simulated typing artifacts by injecting keyboard proximity typos (e.g., "o" for "p") using a synthetic typo generator based on the \textsc{guychuk} dataset distributions \cite{guychuk2024typos}.
\end{itemize}
By using a common set of English source texts for the human modes, we ensured that the differences observed were primarily due to the input 'mode' rather than topical content.

\subsection{Experiment 1: Zero-Shot Classification}
\label{subsec:exp1}

{\bf Goal.} To evaluate whether a state-of-the-art LLM can explicitly identify the origin mode based on text artifacts.

{\bf Setup.} We used the \gptmini model via API with temperature set to 0 to ensure deterministic results.
The model was prompted to act as a forensic linguist and classify each text into one of three categories: \aigen, \dictated, or \typed.
The prompt provided brief definitions of each mode but included no few-shot examples, testing the model's emergent reasoning about text artifacts.

{\bf Metrics.} We recorded Accuracy, Precision, Recall, and F1-Score for each mode.

\subsection{Experiment 2: Latent Embedding Analysis}
\label{subsec:exp2}

{\bf Goal.} To determine if origin mode is inherently captured as a separable latent variable in standard text embedding spaces.

{\bf Setup.} We extracted 384-dimensional sentence embeddings for all 150 samples using the \minilm model.
This model was chosen for its balance of performance and efficiency.
We analyzed the latent space using Principal Component Analysis (PCA) to identify the primary axes of variance.

{\bf Metrics.} We used the Silhouette Score to measure the separation of clusters in the high-dimensional space.
A score of 1 indicates perfect separation, 0 indicates overlapping clusters, and -1 indicates incorrect clustering.
We also calculated the Explained Variance Ratio for the first two principal components.
